(1) 令 \(z\in[0,h]\)。薄膜低Reynolds数的径向动量方程为 \(\mu u_{r,zz}=p_r\)，两盘径向无滑移条件给
\[
u_r(r,z)=\frac{p_r}{2\mu}z(z-h),\qquad
q_r=\int_0^h u_r\,dz=-\frac{h^3}{12\mu}p_r.
\]
由于不可压且盘距随时间变化，厚度积分连续方程为
\[
\dot h+\frac1r\partial_r(rq_r)=0,\qquad
\frac1r\partial_r(rh^3p_r)=12\mu\dot h.
\]
(2) \(h\) 与 \(r\) 无关，轴心压强梯度正则，故 \(p_r=6\mu\dot h\,r/h^3\)。以 \(p(a)=p_a\) 积分得到
\[
p(r)-p_a=\frac{3\mu\dot h}{h^3}(r^2-a^2)
=\frac{3\mu(-\dot h)}{h^3}(a^2-r^2).
\]
中心超压为 \(3\mu(-\dot h)a^2/h^3\)，总法向液压力
\[
F=\int_0^a2\pi r(p-p_a)\,dr
=-\frac{3\pi\mu a^4}{2h^3}\dot h>0.
\]
(3) 忽略上盘惯性，\(F=F_0\)，于是
\[
\dot h=-\frac{2F_0}{3\pi\mu a^4}h^3,\quad
\frac1{h(t)^2}=\frac1{h_0^2}+\frac{4F_0t}{3\pi\mu a^4},
\]
故 \(h(t)=h_0[1+4F_0h_0^2t/(3\pi\mu a^4)]^{-1/2}\)；达到 \(h_1\) 的时间为
\[
t_1=\frac{3\pi\mu a^4}{4F_0}(h_1^{-2}-h_0^{-2})>0.
\]
形式上只有在 \(t\to\infty\) 才接近零厚度。
(4) 剪切率 \(u_{r,z}=p_r(2z-h)/(2\mu)\)。体耗散
\[
\mathcal D=\int_0^a 2\pi r\,dr\int_0^h\mu u_{r,z}^2\,dz
=\int_0^a2\pi r\,dr\frac{h^3p_r^2}{12\mu}
=\frac{3\pi\mu a^4\dot h^2}{2h^3}=F(-\dot h).
\]
单位 \([\mu a^4\dot h/h^3]=\mathrm N\)；排液无限耗时不是微观接触不可发生的真实断言，薄膜润滑、无滑移和连续介质在分子厚度失效。
